% conclusions.tex (v15, 2026-10-21): purpose, evidence, implication, in the order of the Introduction.
\section{Conclusions}\label{sec:conclusions}
ALT observations are few and clustered, and the models that extend them to unobserved cells and years face a trade-off between the transferability of a one-coefficient physical model and the accuracy of machine learning at its training sites. This study shows that the trade-off can be avoided by anchoring the learner to the physics. Evaluated on \nncyAll{} site-year observations of three regions with the climate of each observation year, a gradient-boosting model of the residual of a locally recalibrated Stefan model reduced the Stefan error by \ndrefLeblkRoabs{} cm in the Lena Delta and \ndrefCablkRoabs{} cm in Canada under spatial-block validation, needed only 25--100 local labels for a resolved gain, and stayed within \ndrefLeregRoabs{} cm of the Stefan model when its region was withheld from training, whereas direct machine learning exceeded the Stefan error by \ndrefLeregDzabs{} cm in that situation. Scored on the same sites and years, the published ML products were as accurate as this model only near their own training sites and did not represent the thin active layers of the Lena Delta, where the physics-residual model was more accurate than all five products. The practical consequences are a year-labelled 1~km ALT map of each region for 2011--2024 with an area of applicability and prediction intervals, an evaluation protocol that makes the accuracy of models and products comparable at stated sites and years, and a rule for new regions: start from the physics, add a learner where local labels show that one coefficient misrepresents the sites, and expect a gain from a few dozen labels rather than from thousands.
